
\chapter{Supplementary analyses for on-device memory}
\label{app:on_device_statistics}

This appendix keeps the main argument of Chapter~\ref{ch:on_device_memory} separate from
the details needed to reproduce it. The chapter uses two formal statistical results:
the workstation-to-phone agreement rule and the omnibus question-family test. The
remaining material below describes replicate variation, run-length sensitivity,
thermal peaks and analytical routing scenarios. These secondary calculations are not
additional primary claims.

\section{Cross-device agreement}
\label{app:ch5_parity}

The parity cell used the same 100 stable question identifiers on both devices. Both
runs used Gemma 4 E2B with Vector RAG, \texttt{top\_k=10},
\texttt{num\_ctx=4096}, prompt style v3 and the Nomic embedding prefix. The workstation
used Ollama, while the phone used the quantized llama.cpp path.

\begin{table}[htbp]
  \centering
  \small
  \begin{tabular}{lr}
    \toprule
    Quantity & Result \\
    \midrule
    Paired questions & 100 \\
    Matching correctness labels & 83 (83.0\%) \\
    Required agreement & at least 90\% \\
    Cohen's $\kappa$ & 0.621 \\
    Workstation-only correct & 5 \\
    Phone-only correct & 12 \\
    Exact McNemar $p$ & 0.1435 \\
    Exact answer-string matches & 40 (40.0\%) \\
    \bottomrule
  \end{tabular}
  \caption{Components of the predefined workstation-to-phone agreement check. The rule
  required at least 90\% label agreement and McNemar $p>0.05$. It failed on agreement.}
  \label{tab:app_ch5_parity}
\end{table}

Agreement and McNemar's test serve different purposes. The first asks whether the two
environments label the same questions the same way. The second uses only the 17
discordant pairs and asks whether their direction is imbalanced. The observed $p=0.1435$
does not demonstrate equivalence. It only fails to establish a directional shift. The
83\% agreement therefore remains the operative result, and all Chapter 5 accuracy
numerators come from phone outputs. Figure~\ref{fig:app_parity_gate} illustrates the
predefined parity threshold, measured label agreement, and the paired $2\times 2$ decision matrix.

\begin{figure}[htbp]
  \centering
  \includegraphics[width=\textwidth]{results/thesis_visuals/fig15_parity_gate.pdf}
  \caption{Predefined workstation-to-phone parity gate diagnostic ($n=100$ paired LoCoMo questions).
  Left: required threshold (90\%) vs measured label agreement (83.0\%, Wilson 95\% CI [74.5\%, 89.0\%])
  and exact string match (40.0\%). Right: $2\times 2$ concordance matrix, where ``Box'' is the workstation. The phone scored 37\% and the workstation
  30\%, and the two runtimes worded many answers differently. Under the predefined rule the gate failed,
  so phone accuracy is always measured on the phone.}
  \label{fig:app_parity_gate}
\end{figure}

The result can be regenerated with:

\begin{verbatim}
python parity_check.py \
  results/lap_locomo100_e2b_vr_parity \
  results_phone/phone_locomo100_e2b_vectorrag_n100
\end{verbatim}

The saved report is \texttt{results/parity\_check\_2026-08-04.md}.

\section{Question-family randomization test}
\label{app:ch5_question_family_test}

Chapter~\ref{ch:memory_evaluation} reports 50 LoCoMo questions from each of four
families in every primary Gemma cell. To test the visible category separation without
treating the 200 answers as independent, the analysis works within the ten source
conversations.

For a given reader-memory cell, the observed statistic is Pearson's chi-square from the
$4\times2$ table of question family by correct or incorrect outcome. Under the null,
correctness is exchangeable across family labels within a conversation. Each of 100,000
randomizations shuffles the observed binary labels only within their source
conversation. This keeps the conversation's total number correct and its category
composition fixed. The randomization $p$-value is the proportion of shuffled statistics
at least as large as the observed statistic, with the standard plus-one correction.
Holm's procedure adjusts the eight cell-level $p$-values as one family.
The test was added after the category decomposition was inspected. It is a supporting
analysis, not a member of the primary strategy-comparison family.

\begin{table}[htbp]
  \centering
  \small
  \setlength{\tabcolsep}{4pt}
  \begin{tabular}{llrrrrrr}
    \toprule
    Reader & Memory & C1 & C2 & C3 & C4 & $\chi^2(3)$ & Holm $p$ \\
    \midrule
    E2B & Full context    & 5 & 0 & 10 & 0 & 19.82 & 0.00018 \\
    E2B & Graph triples   & 6 & 6 & 13 & 18 & 12.18 & 0.00548 \\
    E2B & Graph + sources & 8 & 11 & 17 & 30 & 25.78 & $<0.0001$ \\
    E2B & Vector RAG      & 16 & 9 & 16 & 31 & 22.40 & 0.00015 \\
    E4B & Full context    & 6 & 0 & 15 & 0 & 32.08 & $<0.0001$ \\
    E4B & Graph triples   & 6 & 9 & 16 & 24 & 19.34 & 0.00050 \\
    E4B & Graph + sources & 8 & 8 & 18 & 31 & 32.52 & $<0.0001$ \\
    E4B & Vector RAG      & 17 & 10 & 21 & 33 & 23.14 & 0.00015 \\
    \bottomrule
  \end{tabular}
  \caption{Omnibus question-family tests for the eight primary Gemma cells. C1 to C4
  are the numbers correct out of 50 for multi-hop, temporal, open-domain and single-hop
  respectively. The chi-square value is the observed separation statistic. Each
  $p$-value comes from conversation-stratified randomization and is Holm-adjusted across
  all eight rows.}
  \label{tab:app_ch5_question_family}
\end{table}

Every adjusted value is below 0.05. The omnibus test establishes that correctness is
associated with question family within each row. It does not identify which pairs of
families differ and does not turn the adjacent percentages in Chapter~\ref{ch:memory_evaluation} into pairwise findings. The script and source table are:

\begin{verbatim}
python question_family_stats.py
results/question_family_stats.csv
\end{verbatim}

\section{Phone energy replication}
\label{app:ch5_energy_replication}

Table~\ref{tab:app_ch5_energy_replicates} gives the repeated sitting values that are
collapsed to means in Chapter~\ref{ch:on_device_memory}. The standard deviation describes
variation between complete sittings. It is not a per-query error bar.

\begin{table}[htbp]
  \centering
  \small
  \setlength{\tabcolsep}{4pt}
  \begin{tabular}{llrrrrr}
    \toprule
    Reader & Memory & Sitting 1 & Sitting 2 & Sitting 3 & Mean & SD \\
    \midrule
    Qwen3.5 0.8B & Vector RAG   & 248.5 & 323.9 & 345.1 & 305.8 & 50.8 \\
    Qwen3.5 2B   & Vector RAG   & 437.1 & 420.6 & 440.7 & 432.8 & 10.7 \\
    Qwen3.5 4B   & Vector RAG   & 566.7 & 594.3 & 546.1 & 569.0 & 24.2 \\
    Gemma 4 E2B  & Vector RAG   & 653.8 & 473.4 & 525.0 & 550.7 & 92.9 \\
    Gemma 4 E2B  & Full context & 1515.6 & 1338.7 & 1331.9 & 1395.4 & 104.1 \\
    Gemma 4 E4B  & Vector RAG   & 770.9 & 748.3 & 693.3 & 737.5 & 39.9 \\
    \bottomrule
  \end{tabular}
  \caption{Voltage-weighted joules per scored answer in cells with three complete
  sittings. The other full-context cells have one complete sitting.}
  \label{tab:app_ch5_energy_replicates}
\end{table}

Gemma E2B is the only strategy pair with three complete sittings on both sides. Its
within-sitting Vector/full-context correct-per-kilojoule ratios are 23.18, 26.86 and
25.37. Their mean is 25.14 and their standard deviation is 1.85. The second
full-context sitting completed 19 answers, so its accuracy is calculated as $1/19$.
The ratio generator now reads the exact scored labels rather than reusing the first
sitting's 5\% accuracy. This is why the paired mean is not equal to the quotient of the
rounded energy means in the main tables.

Repeated sittings reduce the uncertainty in energy, not in accuracy: greedy decoding gives the
same answers in every sitting, so the accuracy of each cell still rests on 20 questions. A
question-level sensitivity analysis resampled the 20 observed paired questions with replacement
20,000 times, with energy per answer held at its measured mean. The
2.5th percentile of the vector-to-full-context ratio is 2.6$\times$ for Qwen3.5 2B, 5.7$\times$ for
Qwen3.5 4B, 7.6$\times$ for Gemma E2B and 8.8$\times$ for Gemma E4B; for Qwen3.5 0.8B full context
has no correct answer, so no finite ratio is reported. The range 2.6$\times$ to 8.8$\times$ summarizes
four separate lower percentiles, not one interval for the whole experiment. These percentiles are
conditional on fixed energy: they exclude between-sitting energy variation and charge-counter bias,
and the question-level resampling does not account for dependence within a conversation.
In all five readers, every question that full context answered correctly was also answered correctly
by vector retrieval. This describes answer quality on the observed sample; it does not validate the
energy measurements. Resampling these pairs cannot introduce full-context-only successes that were
absent from that sample. The percentiles therefore are not guaranteed lower bounds for new workloads.
The calculation is
\texttt{python phone\_ratio\_bootstrap.py}, which writes \texttt{results/phone\_ratio\_bootstrap.json}.

Accuracy intervals displayed in Figure~\ref{fig:ch5_phone_efficiency} are Wilson 95\%
intervals for 20 binary outcomes. They describe uncertainty in each marginal accuracy
estimate. They do not test whether the efficiency ratio differs across reader families.

\section{Run-length observation}
\label{app:ch5_run_length}

The separate E2B Vector-RAG run over 100 phone questions consumed 29.76 kJ, scored 37
of 100 answers correct and used 297.6 J per answer. The three 20-question sittings used
550.7 J per answer on average. The observed ratio is 1.85. Splitting each run into its
memory-building windows and its question stretches (\texttt{golden\_metric.phase\_split})
accounts for most of it. Building the memory for six histories took 62\% and 60\% of the
energy in the two 20-question sittings that did not start at full charge, against 32\% in the
100-question run, which built ten histories (its first two, begun while the charge counter was
paused at full charge, are left out of that share), while the questions themselves cost about 198~J per
answer at 100 questions against 205 to 246~J at 20. The lower figure at 100 questions therefore
comes mainly from spreading the memory build over more questions. It is still not a matched
run-length intervention: the 100-question run used the task prefixes, and its question
composition and execution time differ. No main-table value is rescaled by 1.85.

This auxiliary cell is stored in \texttt{results/golden\_metric.json} and traced as
\texttt{app\_ch5\_run\_length} in
\texttt{results/thesis\_visuals/figure\_trace.yaml}. It can be regenerated with
\texttt{python golden\_metric.py}.

\section{Thermal peak detail}
\label{app:ch5_thermal}

The mean of the within-sitting maxima of the hottest thermal zone is higher for Vector RAG in
all five reader pairs. Table~\ref{tab:app_ch5_thermal} reports the observation directly.

\begin{table}[htbp]
  \centering
  \small
  \begin{tabular}{lrrr}
    \toprule
    Reader & Full context ($^\circ$C) & Vector RAG ($^\circ$C) & Difference \\
    \midrule
    Qwen3.5 0.8B & 95.80 & 97.53 & 1.73 \\
    Qwen3.5 2B   & 95.40 & 99.13 & 3.73 \\
    Qwen3.5 4B   & 95.80 & 96.33 & 0.53 \\
    Gemma 4 E2B  & 95.67 & 98.87 & 3.20 \\
    Gemma 4 E4B  & 97.00 & 97.93 & 0.93 \\
    \bottomrule
  \end{tabular}
  \caption{Displayed means of per-sitting peaks of the hottest thermal zone. Vector RAG has three
  sittings in every row. Full context has three for E2B and one for the other four
  readers. These values are descriptive and are not battery temperatures.}
  \label{tab:app_ch5_thermal}
\end{table}

The repeated 5-of-5 direction is worth reporting, but unequal replication makes its
uncertainty asymmetric. Peak temperature also ignores exposure duration. The table
therefore does not support a causal claim that retrieval runs hotter.

\section{Long-context and routing sensitivities}
\label{app:ch5_routing_sensitivity}

Figure~\ref{fig:app_ch5_kv_wall} sets the KV-cache size of the two Gemma readers against the
memory the phone had left. The curves are calculated from the layer shapes in the headers of the
phone's own GGUF files, not measured at run time; the horizontal lines are measured. No context
longer than the 8,192-token server context was run on the phone, so the figure describes how the
cache would grow, not how long prompts would perform.

\begin{figure}[htbp]
  \centering
  \includegraphics[width=\textwidth]{results/thesis_visuals/fig16_kv_ram_wall.pdf}
  \caption{Calculated f16 KV-cache size versus context length for Gemma~4 E2B and E4B, using the phone
  GGUF layer shapes. Markers show the reserved 8,192-token cache (57 and 155~MB). Dotted lines
  show measured free RAM at each model's highest system-memory peak (4.4 and 5.0~GB). The curves
  describe cache growth alone; they do not establish total runtime memory or the feasibility of
  longer-context inference.}
  \label{fig:app_ch5_kv_wall}
\end{figure}

KV sharing and 512-token sliding windows leave only three E2B and four E4B layers whose
cache grows with context. These counts agree with the Gemma~4 technical
report~\citep{gemma4}. Its Table~3 gives int8 cache sizes of 0.05 and 0.14~GB at 32,768
tokens, about half the f16 estimates here. The plotted estimates are conservative relative
to those lower-precision footprints; they are not bounds on the runtime's other allocations.

The exploratory routing grid combines the direct per-turn construction anchors with
projected graph query energy and workstation accuracy. It now makes two like-for-like
comparisons. Raw query tokens are compared with raw lifecycle energy, producing different
choices in 38 of 40 history-length and query-count profiles. Query tokens per correct
answer are compared with lifecycle energy per correct answer, producing different
choices in 20 of 40 profiles. These are scenario counts, not statistical tests.
Figure~\ref{fig:app_amortization_frontier} visualizes this dynamic break-even frontier as user query volume scales.

\begin{figure}[htbp]
  \centering
  \includegraphics[width=\textwidth]{results/thesis_visuals/fig17_amortization_frontier.pdf}
  \caption{Illustrative amortization for Gemma~4 E2B over a 419-turn history. (A) Cumulative energy
  including construction; markers show projected crossings with full context. (B) Energy per
  expected correct answer, on a logarithmic scale. Construction and vector/full-context query
  energies are phone measurements; graph query energy is estimated from prompt tokens.
  Workstation accuracy is assumed to transfer despite the failed parity criterion. The phone
  graph-build pilot and workstation graph-accuracy runs also use different configurations. These
  curves describe a scenario, not measured on-device graph efficiency.}
  \label{fig:app_amortization_frontier}
\end{figure}

The full-context and vector query-energy anchors average three sittings. Vector query
energy excludes embedding windows, so construction is counted once. Graph query energy
is scaled from vector query energy using workstation prompt tokens; accuracy uses the
workstation $n=200$ record. The graph build is the early phone pilot, with 58 of 84 failed
batches, whereas the workstation graphs had under 1\% failures. These mismatched inputs
and the failed parity criterion limit the calculation to an illustrative scenario.

Under these assumptions, GraphRAG v3 and v4 cross full context's cumulative energy after
about 17 and 19 queries. Neither overtakes vector retrieval per expected correct answer
over the plotted 1--100 queries: vector construction costs 0.85~kJ against the graph's
21.9~kJ, and its assumed accuracy is higher. These are properties of the projection.

The former 40-of-40 statement is retired because it compared raw prompt tokens with
energy per correct answer. Introducing accuracy normalization changed the result. More
importantly, no complete phone graph-QA cell exists. The only conclusion carried into
the main chapter is that query-token counts omit memory construction. The complete
exploratory calculation is in \texttt{results/routing\_divergence.md} and is regenerated
by \texttt{python routing\_divergence.py}.

\section{Energy and charge conventions}
\label{app:ch5_energy_conventions}

Phone energy integrates each positive charge-counter interval against the mean recorded
battery voltage for that interval. The direct graph construction span is 21.898 kJ for
419 turns, or 52.263 J per turn. Its charge-counter decrease is 1495.305 mAh, which
implies an effective voltage of 4.068 V over that span.

The battery readings provide two consistency checks. In every sitting, the charge-counter
decrease matches the drop in reported battery level, at 47.47~mAh per percentage point, to within
about one percentage point, which is the resolution of the level reading. Over the construction
span, for example, the counter fell 1,495~mAh while the level fell from 91\% to 60\%, or about
1,472~mAh. Integrating the logged instantaneous current against voltage
agrees with the counter energy to within 10\% in every sitting longer than 50 minutes (median
$+2\%$; 22.6 against 21.9~kJ for construction). The instantaneous current is used only for this
check. Termux:API reports it in inconsistent units: values near $-1$ and near $-1{,}000$ describe
the same draw of about one ampere. The logged \texttt{power\_w} column is therefore wrong on most
rows and is not used anywhere in this thesis.

The counter has two limits. It reports in steps of 0.1\% of capacity (4.747~mAh, about 68~J at
4~V), updated about every 30~s, which is negligible over a whole sitting. It also holds its full
value while the battery reads 100\%. Eleven sittings started at the highest possible charge,
100\%. In nine of them the counter stayed flat for 161 to 380~s while
the logged current averaged 0.9 to 1.5~A, and it did not catch up once it began to fall, so that
energy is missing.
Estimated from the logged current over the flat interval, the shortfall is about 2 to 5\% for the
long sittings: E2B full context in its first and third sittings, Qwen3.5 4B full context, and the two
$n=100$ runs. It is about 15 to 25\% for the short ones: E2B vector retrieval in its second sitting,
Qwen3.5 0.8B vector retrieval in its first sitting, Qwen3.5 0.8B full context, and the $k=2$ run. The
reported values are not corrected, because the estimate rests on the inconsistently scaled
instantaneous current. With the estimate applied, the E2B vector-to-full-context ratio would move
from 25.1$\times$ to about 24.5$\times$, and the Qwen3.5 0.8B energies would move the most. No
direction reported in Chapter~\ref{ch:on_device_memory} changes.

Correct answers per full-charge equivalent use measured mAh per answer and the phone's
5,000 mAh capacity. They do not divide by a fixed energy-per-charge constant. The earlier
3.88 V routing convention and its 73\% graph projection are therefore not used in the
chapter. This separation keeps the energy and charge calculations internally consistent.
